\chapter{Conclusiones}
\label{chap:conclusiones}

\section{Respuesta al objetivo}
El análisis de cinco incidentes ocurridos desde septiembre
de 2026 muestra que la relevancia de un suceso de
ciberseguridad depende tanto del fallo técnico como
del contexto en el que se explota. Sus consecuencias
abarcan el acceso a redes internas, la exposición
de información personal, la interrupción de servicios
y la pérdida de bienes físicos.

PHP2Shell e Ivanti ilustran cómo el encadenamiento
de vulnerabilidades puede ampliar los permisos
y el alcance de una intrusión. Moldavia y las entidades
financieras surcoreanas evidencian la importancia
de controlar el acceso a datos desde aplicaciones
y servicios periféricos. Los supermercados franceses
muestran que la seguridad de los sistemas conectados
también protege procesos materiales y la conservación
de alimentos.

Los casos permiten concluir que actualizar es necesario,
pero no suficiente para gestionar un incidente.
La corrección del punto de entrada debe acompañarse
de una investigación del alcance, la recuperación
de los sistemas y medidas frente a las consecuencias
que persisten, como credenciales sustraídas
o datos personales ya extraídos.

\section{Recomendaciones y reflexión final}
A partir de la comparación se proponen tres prioridades.
En primer lugar, las organizaciones deben conocer
sus activos y reducir la exposición de interfaces
administrativas, aplicaciones auxiliares y equipos
sin soporte. En segundo lugar, deben aplicar
actualizaciones y verificar los permisos efectivos
sobre funciones y registros, evitando que un acceso
limitado permita consultar información ajena
o alcanzar sistemas adicionales. En tercer lugar,
la respuesta debe integrar evidencias técnicas,
protección de los afectados y comunicación clara
sobre lo confirmado y lo todavía incierto.

El empleo de IA observado en la campaña surcoreana
añade capacidad de automatización, pero no sustituye
la necesidad de corregir las debilidades explotadas.
Su presencia debe analizarse mediante evidencias,
sin equiparar automáticamente asistencia,
automatización y autonomía completa.

Finalmente, la investigación ha mostrado que
el rigor también consiste en reconocer los límites
de las fuentes. Un CVSS no mide el daño social,
un sistema expuesto no equivale a una víctima
y una hipótesis compatible no demuestra la causa
de un ataque. Las conclusiones se limitan a los
cinco casos estudiados y a la información pública
consultada hasta el 8 de noviembre de 2026;
no constituyen una descripción exhaustiva
del panorama de amenazas.